\subsection{Initial ramification index}

DEFINITION 4. Let G be a totally ordered commutative group and H a subgroup of G of finite index. The number of major subsets of G consisting of strictly positive elements and containing all the elements \textgreater0 of H is called the initial index of H in G and denoted by \(\varepsilon(\mathbf{G}, \mathbf{H})\) .

This initial index is a natural number by virtue of the following proposition:

PROPOSITION 3. Under the hypotheses of Definition 4, if the set of strictly positive elements of G has no least element, then \(\varepsilon(G, H) = 1\) for all H. If there exists a least element \(>0\) of G and G' denotes the subgroup it generates, then

\[
\varepsilon (\mathrm{G}, \mathrm{H}) = (\mathrm{G} ^ {\prime}: (\mathrm{G} ^ {\prime} \cap \mathrm{H})).
\]

In the first case, let x be an element \textgreater0 in G. The set of \(y \in G\) such that 0 \textless{} y \textless{} x is infinite and hence there exist two elements of this set which are distinct and congruent modulo H; their difference is an element z of H such that 0 \textless{} z \textless{} x. Hence every major subset which contains all the strictly positive elements of H contains x and hence all the elements \textgreater0 of G.

In the second case, let \(x\) be the least element \(>0\) of \(G\) and let \(n\) be the least integer \(>0\) such that \(nx \in H\) . Clearly \(n = (G': (G' \cap H))\) . On the other hand, writing \(M(y)\) for the set of \(z \in G\) such that \(y \leqslant z\) ( \(y \in G\) ), it is immediately seen that the major sets of Definition 4 are just \(M(x), M(2x), \ldots, M(nx)\) .

COROLLARY. The initial index \(\varepsilon(\mathbf{G},\mathbf{H})\) divides the index (G: H) and is equal to it if G is isomorphic to Z.

In particular, \(\varepsilon(G, H) \leqslant (G: H)\) .

DEFINITION 5. Let K be a field, L a finite extension of K, w a valuation on L, v its restriction to K and \(\Gamma_{w}\) and \(\Gamma_{v}\) their order groups. The initial index of \(\Gamma_{v}\) in \(\Gamma_{w}\) is called the initial ramification index of w with respect to v (or with respect to K) and denoted by \(\varepsilon(w/v)\) .

From the above corollary, \(\varepsilon(w/v)\) divides \(e(w/v)\) with equality in the case of a discrete valuation.

PROPOSITION 4. Under the hypotheses of Definition 5, let A and m (resp. A' and m') be the ring and ideal of the valuation v (resp. w). Then

\[
[ \mathrm{A} ^ {\prime} / \mathrm{mA} ^ {\prime}: \mathrm{A} / \mathrm{m} ] = \varepsilon (w / v) f (w / v).
\]

The ideals of A' containing mA' and distinct from A' correspond to the major subsets of \(\Gamma_{w}\) consisting of elements \textgreater0 and containing the elements \textgreater0 of \(\Gamma_{v}\) ( \(\S3\) , no. 5, Corollary to Proposition 7). They are therefore equal in number to \(\varepsilon(w/v)\) and, as they form a totally ordered set under inclusion, this number is equal to the length of the quotient ring A'/mA'. Now a module of length 1 over A' is a 1-dimensional vector space over A'/m' and hence a module of length \(f(w/v)\) over A; hence, as A'/mA' is of length \(\varepsilon(w/v)\) over A', it is of length \(\varepsilon(w/v)f(w/v)\) over A, that is over A/m.

